\uuid{59DU}
\exo7id{4236}
\titre{exo7 4236}
\auteur{quercia}
\organisation{exo7}
\datecreate{2010-03-12}
\isIndication{false}
\isCorrection{true}
\chapitre{Calcul d'intégrales}
\sousChapitre{Théorie}
\cpge{mpsi-mp2i.primitives-equations-differentielles.a, pcsi.techniques-calcul-primitives.a, ptsi.techniques-calcul-primitives.a, mpsi-mp2i.integration.c, pcsi.integration.b, ptsi.integration.b, bcpst1.derivees-primitives-integrales.d, bcpst1.integration-segment.d}
\module{Analyse}
\niveau{L1}
\difficulte{}

\contenu{
\question{
Calculer $I =  \int_{t=0}^\pi \ \frac t{1+\sin t}\,d t$.
}
\reponse{
$u = \pi-t  \Rightarrow  I = \frac \pi2 \int_{t=0}^\pi \ \frac 1{1+\sin t}\,d t
= \frac \pi2  \int_{t=-\pi/2}{\pi/2} \ \frac 1{1+\cos t}\,d t = \pi$.
}
}
